\chapter{Uso de herramientas de inteligencia artificial}\label{uso-ia}
Durante la elaboración de este Trabajo de Fin de Grado se han utilizado herramientas de inteligencia artificial generativa como apoyo auxiliar en tareas de redacción, revisión, organización y documentación. Su uso se ha limitado a facilitar el proceso de trabajo, sin sustituir el criterio técnico del autor ni la validación final de los contenidos incluidos en la memoria.

\section{Objetivo del uso}

El uso de estas herramientas ha tenido como objetivo mejorar la claridad de la documentación, contrastar la coherencia entre apartados y acelerar tareas mecánicas asociadas a la preparación de la memoria. En particular, se han empleado como apoyo para estructurar explicaciones, revisar formulaciones, detectar posibles incoherencias y generar borradores iniciales que posteriormente han sido corregidos y adaptados.

\section{Tareas realizadas con apoyo de IA}

Las herramientas de inteligencia artificial se han utilizado principalmente en las siguientes tareas:

\begin{itemize}
    \item Revisión y mejora de fragmentos de texto de la memoria.
    \item Apoyo en la organización de secciones y anexos del documento.
    \item Comprobación de coherencia entre requisitos, casos de uso, diagramas y explicaciones.
    \item Generación de propuestas iniciales de diagramas y material auxiliar, revisadas y modificadas posteriormente por el autor.
    \item Apoyo en tareas de compilación, exportación y mantenimiento de la documentación del proyecto.
\end{itemize}

\section{Supervisión y responsabilidad del autor}

Todos los contenidos incorporados al documento han sido revisados, corregidos y validados por el autor. Las decisiones sobre el alcance del sistema, la selección de requisitos, la interpretación de los diagramas, la redacción final y la estructura de la memoria corresponden al autor del trabajo.

La inteligencia artificial no se ha utilizado como fuente primaria de información académica o técnica. Cuando se han incluido definiciones, criterios de notación o referencias externas, se han utilizado fuentes documentales específicas citadas en la memoria. Por tanto, el uso de IA debe entenderse como una herramienta de apoyo al proceso de elaboración, no como sustituto de la documentación formal ni del trabajo propio realizado.

